\chapter{Approach} \label{approach}
This chapter presents our experimental approach, which includes the evaluation metrics and models for both classification and detection tasks. We will also discuss the specific GAN-generated architectures used and provide an overview of the overall pipeline for the experiments.

\section{Metrics}
In this session, we will discuss the evaluation metrics used in this thesis and their significance. The evaluation metrics are categorized into classification metrics and detection metrics, which are specific to their respective tasks. The classification metrics include accuracy, recall, precision, and F1-score, while the detection metrics include Intersection over Union (IoU), mean average precision (mAP), precision, class loss, object loss, and classification loss. We will provide a brief explanation and overview of each metric.
\subsection{Classfication Metrics}
In classification tasks, evaluation metrics are used to measure the accuracy of a model in predicting the correct class label for an input instance. Some of the most important metrics in classification tasks used in this thesis are accuracy, precision, recall, F1-score and their top-N metrics:
\subsubsection{Accuracy, precision, recall and F1-score}
\begin{itemize}
    \item Accuracy is a commonly used metric that measures the overall correctness of the predictions made by a model, expressed as the percentage of correctly classified samples over the total number of samples. 
    \item Precision, also known as positive predictive value, measures the proportion of true positives that are correctly identified by the model over the total number of positive predictions. This metric is important in cases where false positives can have significant consequences, such as in the case of misidentifying an important class.
    \item Recall, also known as sensitivity or true positive rate, measures the proportion of true positives that are correctly identified by the model over the total number of actual positives. This metric is particularly important in cases where false negatives can have severe consequences, such as in the detection of rare species.
    \item The F1 score is a weighted average of precision and recall, and provides a more balanced measure of the model's performance. It is calculated as the harmonic mean of precision and recall. It is calculated as:
    \begin{equation}
        \frac{2 \cdot (precision \cdot recall)}{(precision + recall)}
    \end{equation}
\end{itemize}

While precision and recall are both important metrics for different reasons - precision measures the model's ability to correctly identify true positives, while recall measures the model's ability to identify all positive instances -  F1 score is often considered the most relevant metric for evaluating a machine learning model. By providing a balanced measure of the model's precision and recall,  F1-score offers a more comprehensive evaluation of a model's performance. This is especially useful for imbalanced datasets, such as ours, where precision and recall alone may not accurately represent the model's performance. Therefore, in our thesis, we will use F1-score, beside accuracy, as the primary metric to compare significant outcomes.

\subsubsection{Top N Metrics}

When evaluating a classification model with a high number of classes, such as in our case with more than 100 classes, top-N metrics such as top-N accuracy or top-N F1 score can be useful. Top-N metrics are a variation of the traditional metrics that consider a prediction to be correct if the true class label is among the N highest-ranked predicted class labels. This is particularly useful when a model can provide a ranked list of predictions, as it allows for partial credit to be given for correct predictions that may not be the top-ranked prediction. The specific value of N can be chosen based on the problem domain and use case, typically a small value such as 1 or 5. In our case, we will use top-1 and top-3 metrics.

To illustrate the concepts of top-N metrics, let's use a sample of class ``herring gull''  or ``Silbermöwe'', as illustrated in Figure \ref{fig:twobirds}a).  For this given test example, EfficientNet produces a ranked list of predicted labels, along with their associated probabilities as follows: 

\begin{itemize}
    \item Herring Gull (Silbermöwe): 60\%
    \item Common Gull (Sturmmöwe): 20\%
    \item Buoy  or mooring ring (Boje oder Bojenkette): 10\%
    \item Mallard Duck ( Stockente): 10\%
\end{itemize}

If the prediction in this example was Common Gull, it would be considered a false prediction in top-1 accuracy, but true in top-3 accuracy. This example also highlights the importance of using top-N metrics, as our dataset contains over 100 classes, and some bird classes, such as Herring Gull and Common Gull in Figure \ref{fig:twobirds}, have very similar appearances. Thus, we can still give credit to the models when their predictions are close to the actual labels.

\begin{figure}
\hfill
\subfigure[Herring Gull]{\includegraphics[width=5cm]{src/background/Silbermoewe.png}}
\hfill
\subfigure[Common Gull]{\includegraphics[width=5cm]{src/background/sturmmoewe.jpg}}
\hfill
\caption{Illustration of (a) Herring Gull (Silbermöwe) and (b) Common Gull (Sturmmöwe) to illustrate similar-looking classes}
\label{fig:twobirds}
\end{figure}

\subsection{Detection Metrics}
In detection tasks, it is crucial to evaluate the accuracy of a model in both localizing and classifying objects in an image. Therefore, besides the metrics used in classification tasks such as precision or recall, there are metrics specifically designed to measure the accuracy of predicted localization. This section will introduce some of the important statistics, including Intersection over Union (IoU), Mean Average Precision (mAP), and various loss metrics such as class loss, object loss, and classification loss used during training.
\subsubsection{Intersection over Union (IoU)}
IoU is a measure of the overlap between predicted bounding boxes and ground-truth bounding boxes. IoU is calculated as the ratio of the area of intersection between the predicted and ground-truth bounding boxes to the area of their union. It ranges from 0 to 1, where 1 indicates perfect overlap and 0 indicates no overlap. IoU is commonly used as a threshold for determining true positive or false positive detections, with 0.5 being a common threshold. Figure \ref{fig:iou} illustrates the intersection over union calculations.

\begin{figure}
	 \centerline{\includegraphics[width=0.8\textwidth]{src/background/iou.png}}
	 {\caption{Illustration of the Intersection over Union  calculation \cite{map}} \label{fig:iou}}
\end{figure}

\subsubsection{Mean Average Precision (mAP)} \label{mAP}
The average precision (AP) is a measure of the accuracy of an object detection model in detecting objects of a particular class. AP is calculated by computing the precision and recall for a range of IoU thresholds, and then calculating the area under the precision-recall curve (PR curve).

The following steps describe how to calculate the AP at a certain IoU:
\begin{itemize}
    \item For each object class, calculate the precision and recall.
    \item Plot the precision-recall curve by plotting the precision values against the corresponding recall values at each IoU threshold.
    \item Calculate the area under the precision-recall curve. The resulting value is the AP for that object class.
\end{itemize}

AP is a useful metric for measuring the accuracy of an object detection model for a specific class, but its calculation is dependent on the IoU threshold, which can be arbitrary and varies for different tasks. To address this ambiguity, mean average precision (mAP) was introduced as a comprehensive metric that evaluates the model's accuracy across multiple object classes and IoU thresholds.

Calculating mAP is a simple process that builds upon the calculation of AP. Firstly, a set of thresholds is considered for AP calculation. In this thesis, we use thresholds ranging from 0.5 to 0.95 with a step size of 0.05 (mAP@[0.5:0.95]). AP is then calculated across the set of IoU thresholds for each object class, and the average of all AP values is taken to obtain the mAP. mAP@[0.5:0.95] is a widely used metric for evaluating object detection models because it provides a comprehensive evaluation of the model's accuracy across a reasonable range of IoU levels. 

\subsubsection{Loss Metrics}
During training we use box loss, classification loss, and object loss. 

\begin{itemize}
    \item Box loss measures the discrepancy between predicted and ground-truth bounding box coordinates. It encourages the model to better localize objects in an image.
    \item Classification loss measures the discrepancy between predicted and ground-truth class probabilities. It encourages the model to correctly classify objects in an image.
    \item Object loss is a combination of box loss and classification loss, and is used to evaluate the overall accuracy of the model in localizing and classifying objects in an image.
\end{itemize}


    

    
\section{Evaluation Models}
In this section, we provide a brief overview of two popular deep learning models that we use to evaluate our augmentation methods: YOLOv5s for object detection and EfficientNet version B0 for image classification. The evaluation is based on F1-score for classification task and mAP@[0.5:0.95] for detection task. Understanding these models is essential to accurately interpret the results of our evaluation.

\subsection{EfficientNet}
There are numerous classification models available, such as YOLOv8 Classification, ResNet, EfficientNet, and Vision Transformer. Despite their good performance, as demonstrated in benchmarks like ImageNet \cite{benchmark}, we have opted for EfficientNet due to its ease of modification and implementation using the Timm library \cite{wightman2020timm}.

EfficientNet \cite{tan2019efficientnet} is a family of neural network models that achieve state-of-the-art performance on various computer vision tasks. It is a convolutional neural network architecture designed to achieve high accuracy while keeping the number of parameters low. This is achieved through a technique called ``compound scaling'' where the network's depth, width, and resolution are scaled in a balanced way, but independently based on a fixed scaling coefficient. This approach is different from traditional neural network scaling, where the dimensions are increased simultaneously by a fixed factor, which can lead to overfitting, computational inefficiency, and decreased accuracy.

This approach allows EfficientNet to achieve better performance than previous models, such as ResNet and Inception, while requiring fewer parameters. This makes EfficientNet an attractive choice for our classification task, as we are dealing with a large dataset with more than 100 classes and limited computational resources. Another advantage of using EfficientNet is that it has been pre-trained on a large-scale dataset such as ImageNet, which can help improve the performance of our model on our specific dataset. Figure \ref{fig:efficientnet} provides a more comprehensive understanding of the performance of EfficientNet between different models \cite{tan2019efficientnet}.

\begin{figure}[!h]
	 \centerline{\includegraphics[width=0.6\textwidth]{src/background/efficientNet.png}}
	 {\caption{Model Size vs. Accuracy Comparison. EfficientNet-B0 is the baseline network developed by AutoML MNAS, while Efficient-B1 to B7 are obtained by scaling up the baseline network. In particular, EfficientNet-B7 achieves new state-of-the-art 84.4\% top-1 / 97.1\% top-5 accuracy, while being 8.4x smaller than the best existing CNN.  \cite{tan2019efficientnet}} \label{fig:efficientnet}}
\end{figure}

EfficientNet models come in several sizes, from EfficientNet B0 (the smallest) to EfficientNet B7 (the largest). The larger models have more layers and more parameters, and are capable of achieving higher accuracy, but require more computational resources to train and run. This thesis uses pretrained EfficientNet with version B0. Although the performance may be affected by using version B0, it aligns with the goal of our thesis to compare various augmentations, making it acceptable. By using version B0, which has fewer parameters, we can train the model faster with different configurations, resulting in quicker analysis. Figure \ref{fig:efficientnetb0} shows the architecture of our baseline network EfficientNet-B0.

\begin{figure}[!h]
	 \centerline{\includegraphics[width=1\textwidth]{src/background/efficientnetb0.png}}
	 {\caption{The architecture for our baseline network EfficientNet-B0 \cite{tan2019efficientnetArticle}} \label{fig:efficientnetb0}}
\end{figure}


\subsection{YOLOv5 (You Only Look Once version 5)}

YOLOv5 (You Only Look Once version 5) \cite{glenn_jocher_2022_7347926}  is a state-of-the-art object detection model developed by Ultralytics. It is inspired by the success of previous versions of YOLO \cite{bochkovskiy2020YOLOv4} and introduces several improvements, including a new backbone architecture and a more efficient training process (mixed precision training \cite{micikevicius2017mixed} and gradient accumulation \cite{hermans2017accumulated}).

Figure \ref{fig:YOLO-general} shows the high-level architecture for an object detector like YOLOv5 which is comprised of three components: backbone, neck, and head \cite{YOLOarticle}. 

\begin{itemize}
    \item The model backbone: a pre-trained network used for feature extraction from the images, which reduces the spatial resolution while increasing the feature (channel) resolution.
    \item The model neck: responsible for extracting feature pyramids, which enables the model to effectively detect objects of varying sizes and scales, leading to better generalization performance.
    \item The model head: responsible for performing the final operations, which involve applying anchor boxes to feature maps and producing the final output. This output includes information such as object classes, objectness scores, and bounding box coordinates.
\end{itemize}

\begin{figure}[!h]
	 \centerline{\includegraphics[width=1\textwidth]{src/background/YOLO-general.png}}
	 {\caption{Single-Stage Detector Architecture\cite{bochkovskiy2020YOLOv4}} \label{fig:YOLO-general}}
\end{figure}

The network architecture of YOLOv5 is shown in Figure \ref{fig:YOLOv5s}, where the backbone architecture is based on CSPNet (Cross Stage Partial Network) \cite{wang2020cspnet}. CSPNet is a novel CNN design that divides a standard convolution layer into two parallel streams, reducing computation costs and enhancing the model's ability to capture complex features. The Neck of YOLOv5 incorporates Spatial Pyramid Pooling (SPP) and Path Aggregation Network (PANet). SPP enables object recognition of different sizes without image resizing or cropping by dividing the image into a fixed number of regions and pooling features from each region. PANet utilizes a bottom-up pathway to extract features from the input image and a top-down pathway to generate a feature map for object detection, effectively aggregating multi-scale features to capture objects at different scales and resolutions. The head of YOLOv5 is composed from three convolution layers that predicts the location of the bounding boxes (x, y, height, width), the scores and the objects classes.

\begin{figure}[!h]
	 \centerline{\includegraphics[width=1\textwidth]{src/background/YOLOv5s.png}}
	 {\caption{Network architecture for YOLOv5e \cite{xu2021forest}} \label{fig:YOLOv5s}}
\end{figure}

There are different sizes available in YOLOv5, ranging from YOLOv5n with 1.9 million trainable parameters to YOLOv5x6 with 140.7 million trainable parameters. As the model size increases, it becomes more capable of achieving higher mAP \cite{glenn_jocher_2022_7347926}. However, to ensure our computational resources are not exceeded, we will use YOLOv5s with 7.2 million trainable parameters.
\section{Training Pipeline}
This section presents a detailed description of the experimental set-up and training pipeline used in our thesis.

The experiments were performed on a computer with an AMD Ryzen Threadripper 2970WX 24-Core Processor and 188 Gb RAM. The deep learning models Yolov5s and GANs were trained on an NVIDIA GeForce RTX 3090 graphics card, while the EfficientNet model was trained on an NVIDIA GeForce RTX 2080 Ti. The software environment included Python 3.10.8, Pytorch, and the Albumentations library for data augmentation.
\subsection{Pipeline for basic augmentations}
Figure \ref{fig:basic_pipieline} illustrates the experiment pipeline for non-GAN approaches. For the basic augmentations, including geometric, color-space, cutout, and mixed augmentations, we apply online augmentation during training for the classification task. This involves directly applying the transformations to the images during training and then feeding them to the EfficientNet version B0. However, for the detection task using YOLOv5s format, we perform the augmentations offline on the dataset before feeding it into the model. 

\begin{figure}[!h]
	 \centerline{\includegraphics[width=0.7\textwidth]{src/approach/basic_pipeline.pdf}}
	 {\caption{Experiment process for basic augmentations (Color-space, geometrical, cutout and mixed augmentations)} \label{fig:basic_pipieline}}
\end{figure}

\subsection{Pipeline for GAN-generated augmentation}
Figure \ref{fig:gan_pipieline} shows the pipeline for GAN-generated methods. We focus on improving the edge cases in our dataset. To achieve this, we only train the GAN on classes that have less than 10 samples, which we consider as edge cases. Once the GAN is trained, we use it to generate half the amount of available images for the edge cases and add them to the training set. We then train this extended training set without any additional augmentation and compare the results to the baseline. It is also important to note that we pay special attention to the underrepresented classes in our dataset. This is because we want to ensure that our models perform well on all classes, not just the most frequently occurring ones. By comparing the results of our GAN-generated methods on the edge cases to the baseline, we can see if the model's performance has been significantly improved for these underrepresented classes.

\begin{figure}[!h]
	 \centerline{\includegraphics[width=0.8\textwidth]{src/approach/gan_pipeline.pdf}}
	 {\caption{Experiment process for GAN-generated approach} \label{fig:gan_pipieline}}
\end{figure}

The performance of GANs heavily relies on the specific architecture of the networks used. Both vanilla GAN and Wasserstein GAN have their own unique architectures that can significantly impact their performance. In the following sections, we will delve into the details of the architectures used for vanilla GAN and Wasserstein GAN.

\subsubsection{Vanilla GAN architecture}
\begin{figure}[!h]
	 \centerline{\includegraphics[width=0.9\textwidth]{src/approach/vanilla_architecture.pdf}}
	 {\caption{he architecture of our generator and discriminator in Vanilla GAN } \label{fig:vanilla_architecture}}
\end{figure}

The architecture of our vanilla GAN, depicted in Figure \ref{fig:vanilla_architecture}, uses a concatenated input of the one-hot coded label and noise for the generator. The generator is composed of four dense layers, each followed by LeakyReLU activation layer with a slop of 0.2. Starting with 512 neurons, the output has a size of 196608 at the end, which is then reshaped into (3, 256, 256) after going through tanh() function for normalization to the range of -1 and 1. For the discriminator, the input is flattened and concatenated with the embedded label. The discriminator also has four dense layers, each followed by LeakyReLU activation layer with a slop of 0.2 and Dropout with a probability of 0.3. The output after the sigmoid layer is a single value between 0 and 1, predicting whether the image is real or fake.


\subsubsection{Wasserstein GAN architecture}
The architecture used in our experiments is illustrated in Figure \ref{fig:wgan_architecture}, which is based on a DCGAN structure \cite{radford2015unsupervised}. In the generator, the input noise is concatenated with the labels that are embedded into a one-hot vector, and fed into a ConvTranspose2D block (also known as a 2D Deconvolution block), followed by batch normalization and ReLU activation. We start with 8192 features and halve the depth while doubling the height and width for every subsequent block until we obtain a 256x256 image with 3 channels after 6 blocks, for a model with 256 features. Finally, we use the Tanh() layer to normalize the output within the range of -1 to 1.

The discriminator architecture consists of a dense layer with 65536 units after one-hot encoding for the label input. The output is then reshaped into (1, 256, 256) and concatenated with the image to predict. This concatenated input is given to a block of 4x4 Conv2D layer, followed by an instance normalization layer (except for the first layer) and a LeakyReLU activation layer with a slop of 0.2. At each subsequent Conv2D layer, the height and width are reduced by half, but the depth is doubled. Finally, at the last layer, the input of (8192, 4, 4) is passed through a convolutional layer with a filter size of 4x4 to produce a single-valued output, which gives the label of the class.
\begin{figure}[!h]
	 \centerline{\includegraphics[width=0.8\textwidth]{src/approach/wgan_architecture.pdf}}
	 {\caption{The architecture of our generator and discriminator in WGAN \label{fig:wgan_architecture}}
\end{figure}

\subsection{WGAN-GP architecture}
Our WGAN-GP architecture is inspired by implementation of Jalola \cite{wgangp}. For WGAN-GP, we use a different architecture than WGAN by incorporating Upsampling Residual Blocks instead of ConvTranspose2D layers for the Generator and Downsampling Residual Blocks for the Discriminator. The Generator architecture is depicted in Figure \ref{fig:wgan_gp_gen}. The Upsampling Residual Block contains an UpSampleConv layer after the input, which is used as a shortcut to the output. This is followed by BatchNorm, ReLu layers, and Conv2D layer, which are repeated once more before being concatenated with the shortcut to produce the output. The Upsampling Residual Block has the same effect on the input size as the ConvTranspose2D, halving the depth while doubling the height and width. In the generator, the embedded Label and noise are concatenated and provided as input to a Conv2d layer, which is followed by four Upsampling Residual Blocks before being given to the another chain of BatchNorm, Conv2D and ReLU layers. Finally, the output is normalized using the tanh() layer.

\begin{figure}[!h]
	 \centerline{\includegraphics[width=0.8\textwidth]{src/approach/wgan_gp_gen.pdf}}
	 {\caption{The architecture of our generator in WGAN-GP} \label{fig:wgan_gp_gen}}
\end{figure}

Figure \ref{fig:wgan_gp_dis} illustrates the architecture of the discriminator in WGAN-GP. The Residual Block of the discriminator differs from that of the generator in that it uses DownSampleConv instead of UpSampleConv right after the input. The image is concatenated with the embedded label and passed through a dense layer to reduce the number of neurons. This is then followed by four blocks of Residual Block with Downsampling before being given to a dense layer with an output size of one, which produces the final prediction.

\begin{figure}[!h]
	 \centerline{\includegraphics[width=0.8\textwidth]{src/approach/wgan_gp_dis.pdf}}
	 {\caption{The architecture of our discriminator in WGAN-GP} \label{fig:wgan_gp_dis}}
\end{figure}

% In the context of our thesis, we can use conditional GANs to generate new images of marine wildlife for training our detection and classification models. We can train a GAN on a set of real images of marine wildlife and use it to generate new images that augment the training dataset. For example, we can use a GAN to generate new images of a bird species, such as the black guillemot (gryllteiste), that are underrepresented in our training dataset. However, training a GAN is computationally expensive, and since our dataset has 117 classes, it is not feasible to train a model large enough to memorize all features of each class. Therefore, we will reduce the dimension of the classes to only 10 classes that have few samples. By doing so, we can improve the performance of our marine wildlife detection and classification models.


